\documentclass[10pt]{article}
\usepackage{hyperref}
\usepackage{enumitem}
\usepackage{cmap}
\usepackage[utf8]{inputenc}
\usepackage[T1]{fontenc}
\usepackage{lmodern}
\usepackage[english]{babel}
\usepackage[left=1.05cm,top=0.82cm,right=1.05cm,bottom=0.68cm]{geometry}
\IfFileExists{glyphtounicode.tex}{\input{glyphtounicode}\pdfgentounicode=1}{}
\hypersetup{
    hidelinks,
    pdfauthor={Kevin Bravo},
    pdftitle={Kevin Bravo - Software Developer Mid Senior - 789.mx},
    pdfsubject={CV ATS para la posicion Software Developer Mid Senior en 789.mx},
    pdfkeywords={Software Developer, Backend, TypeScript, JavaScript, Node.js, PostgreSQL, SQL, APIs REST, microservicios, AWS, Java, C Sharp, SQL Server, MongoDB, software empresarial, infraestructura, cloud}
}
\setlength{\parindent}{0pt}
\hyphenpenalty=10000
\exhyphenpenalty=10000
\emergencystretch=2em
\pagestyle{empty}
\setlist[itemize]{leftmargin=*, noitemsep, topsep=0pt, partopsep=0pt, parsep=0pt, label={\textbf{-}}}

\newcommand{\sectiontitle}[1]{%
    \vspace{4pt}
    \begin{center}
        \textbf{#1}
    \end{center}
    \vspace{-5pt}
}

\begin{document}
\raggedright

\begin{center}
    {\LARGE \textbf{Kevin Bravo}}\\
    Software Developer Mid/Senior | Backend y Sistemas Empresariales \\
    \hrulefill
\end{center}

\begin{center}
    \href{mailto:0bravokevin@gmail.com}{0bravokevin@gmail.com}
    | \href{https://kevinbravo.com}{kevinbravo.com}
    | \href{https://github.com/0bkevin}{github.com/0bkevin}
    | \href{https://www.linkedin.com/in/0bkevin/}{linkedin.com/in/0bkevin}
    | Remote
\end{center}

\fontsize{10}{11.3}\selectfont

\sectiontitle{Perfil Profesional}
Desarrollador de software con 5 años de experiencia convirtiendo operaciones complejas en sistemas empresariales claros, mantenibles y medibles. He acompañado productos desde el levantamiento de necesidades y el modelado de datos hasta APIs, interfaces, despliegue, soporte y mejora continua. Trabajo principalmente con TypeScript, JavaScript, Node.js, Python, Django, PostgreSQL y SQL; también cuento con fundamentos transferibles para integrarme a microservicios en C\# o Java y a entornos con SQL Server, MongoDB y AWS. Me interesa el enfoque de 789.mx: entender cómo funciona un negocio, reducir fricción y construir la infraestructura que le permita operar y crecer con mayor control.

\sectiontitle{Experiencia Profesional}

\textbf{Proyecto Educa} \hfill \textit{Remoto, Venezuela} \\
\textbf{Chief Technology Officer} \hfill \textit{Dic 2025 - Actualidad}
\begin{itemize}
    \item \textbf{Traduzco necesidades operativas en arquitectura, prioridades y entregas} para plataformas educativas; coordino colaboradores y stakeholders y respondo por la calidad de releases y la continuidad en producción.
\end{itemize}

\vspace{2pt}
\textbf{Free2Z} \hfill \textit{Remoto, Venezuela} \\
\textbf{Principal Software Engineer} \hfill \textit{Abr 2025 - Abr 2026} \\
\textbf{Open-Source Contributor} \hfill \textit{Jul 2026}
\begin{itemize}
    \item \textbf{Lideré una migración en producción de React a SvelteKit} y desarrollé flujos completos sobre APIs REST en Django, PostgreSQL, Docker, GKE, Cloud SQL, GCS, Pub/Sub y Cloud Build.
    \item \textbf{Fortalecí procesos asíncronos, webhooks y despliegues} con idempotencia, transacciones, reintentos, ACK/NACK, health checks, rollback automático y estados de fallo explícitos.
    \item \textbf{Protegí la continuidad de la plataforma} con migraciones concurrentes de PostgreSQL, pruebas de integración y concurrencia sobre una base real, logging y depuración operativa.
\end{itemize}

\vspace{2pt}
\textbf{U.S. Department of State (a través de AVAA)} \hfill \textit{Remoto, Venezuela} \\
\textbf{Product Engineer | Billingua Talent} \hfill \textit{Nov 2024 - Ago 2025}
\begin{itemize}
    \item \textbf{Construí una plataforma B2B multi-tenant como único ingeniero}, desde el diagnóstico con stakeholders hasta el modelo relacional en PostgreSQL, permisos por roles, APIs, dashboards, despliegue y mantenimiento.
    \item \textbf{La validé con 50 usuarios y 5 empresas} y convertí su retroalimentación en flujos más simples, datos más confiables y mejores reportes para la operación.
\end{itemize}

\vspace{2pt}
\textbf{Freelance} \hfill \textit{Remoto} \\
\textbf{Software Engineer} \hfill \textit{Sep 2024 - Dic 2024}
\begin{itemize}
    \item \textbf{Entregué productos web y automatizaciones a la medida} para distintos sectores, cubriendo descubrimiento, priorización, desarrollo, lanzamiento, iteración y soporte técnico.
\end{itemize}

\vspace{2pt}
\textbf{Asociación Venezolano Americana de Amistad (AVAA)} \hfill \textit{Caracas, Venezuela} \\
\textbf{Junior Product Engineer} \hfill \textit{Oct 2022 - Jul 2024} \\
\textbf{Software Developer Intern} \hfill \textit{Sep 2021 - Ago 2022}
\begin{itemize}
    \item \textbf{Diseñé, desplegué y mantuve SEP}, un sistema multi-tenant con Next.js, PostgreSQL, Prisma, Auth.js y Azure que centralizó actividades, asistencia, voluntariado y reportes de unos 300 becarios en 3 capítulos.
    \item \textbf{Construí SEA para EducationUSA} para integrar historias de estudiantes, servicios, pagos, recibos, actividades, alianzas y reportes; normalicé datos históricos mediante ETL en Python.
    \item \textbf{Mantuve ambos sistemas en producción}, resolviendo defectos, cuidando la calidad de los datos y ajustando procesos con base en la retroalimentación del equipo operativo.
\end{itemize}

\sectiontitle{Voluntariado Relevante}
\textbf{AVAA | Miembro del Comité de Educación} \hfill \textit{Sep 2024 - Actualidad} \\
Administro SEP y mantengo sus flujos de asistencia, voluntariado y reportes para unos 300 becarios; asesoro requisitos, calidad de datos y viabilidad técnica. \\
\textbf{JA Venezuela Alumni | Vicepresidente} \hfill \textit{Mar 2026 - Actualidad} \\
Defino prioridades con la junta y coordino operaciones, comités, alianzas y proyectos de la red nacional de egresados.

\sectiontitle{Educación e Idiomas}
\textbf{Instituto de Estudios Superiores de Administración (IESA)} \hfill \textit{Caracas, Venezuela} \\
Maestría en Gerencia Pública, en curso \hfill \textit{2025 - 2027} \\
\textbf{Idiomas:} Español (Nativo), Inglés (C1 - competencia profesional).

\end{document}
